\documentclass[10pt,a4paper]{amsart}
\usepackage[margin=19mm]{geometry}
\usepackage{amsmath,amssymb,mathtools}
\usepackage[scr=rsfs,cal=euler]{mathalfa}
\usepackage{xfrac}
\usepackage[unicode,hidelinks,bookmarks=false]{hyperref}
\newcommand{\ie}{i.e.\ }
\newcommand{\cf}{cf.\ }
\newcommand{\resp}{respectively\ }
\newcommand{\Subsection}{\subsection}
\providecommand{\nobreakdash}{\nobreak\hbox{-}}
\setlength{\emergencystretch}{2em}
\multlinegap0pt
\usepackage[shortlabels]{enumitem}
\usepackage{amssymb}
\usepackage{mathtools}
\newtheorem{lemma}{Lemma}
\newtheorem{theorem}{Theorem}
\newcommand{\R}{\mathbb R}
\newcommand{\Sk}{\operatorname{Sk}}
\newcommand{\Z}{\mathbb Z}
\newcommand{\arch}{\mathrel{\asymp}}
\newcommand{\st}{\operatorname{st}}
\title{Finite-tower bounds for Skolem functions}
\author{Andreas Weiermann}
\date{Source: arXiv:2609.12740v1 (CC BY 4.0)}
\begin{document}
\maketitle

\noindent\emph{Source and licence.} Finite-tower bounds for Skolem functions. Version arXiv:2609.12740v1 (CC BY 4.0), distributed by arXiv under the Creative Commons Attribution 4.0 International licence, http://creativecommons.org/licenses/by/4.0/, as recorded in the arXiv licence metadata for this exact version and shown on the abstract page https://arxiv.org/abs/2609.12740v1. Full licence notice: \texttt{licenses/arxiv-2609.12740-CC-BY-4.0.txt}.\par\smallskip

\noindent\textit{Editorial scope.} Counted unit: the article's theorem sf:thm:discrete (source0.tex lines 461--470). The article's complete original proof of it is delivered with it (source0.tex lines 471--655, one \texttt{proof} environment of 185 lines). No mathematics is written, repaired or re-derived: the statement and the proof are the article's own lines, in the article's own notation, and no retained line is substituted or re-flowed.  Retained uncounted as the source-local context the proof needs: lines 361--363; lines 376--392; lines 415--422; lines 440--444.  Nothing was omitted from any retained span, so the package carries no omission record.  No retained line contains a citation, so the wrapper carries no bibliography and no bibliography entry.  No retained line contains a figure environment, an image-inclusion command or drawing code, so no image is delivered and the package declares no asset dependency.  Editorial formatting only, in addition: an AMS \texttt{amsart} preamble that restates the article's own declarations and macros used by the retained lines, namely the environments \texttt{lemma}, \texttt{theorem} on separate counters, together with the standard packages those lines need.  The retained environments print in this document as Lemma 1, Theorem 1 in order of appearance, whereas the article prints the same items with the numbering of its own full document.  The complete native source of the article as distributed in the arXiv e-print for this exact version is bundled unchanged as \texttt{sources/210147/source0.tex}.
\begin{lemma}[Bounded Skolem differences]\label{sf:lem:integers}
If $f,g\in\Sk$ and $f-g=O(1)$ in $F$, then $f-g\in\Z$.
\end{lemma}

\begin{lemma}[Scaled comparisons]\label{sf:lem:scaled}
Let $c\in F$, $c\ge1$.
\begin{enumerate}[label=\textup{(\roman*)},leftmargin=2.5em]
\item The relation $u\arch_c v$ defined by $(u/v)^c\arch1$ is
an equivalence relation with convex classes on $F_{>0}$. It implies
$u\arch v$.
\item If $c$ is infinite and $z>0$, then for every $s\in\R$,
\[
 z^c=e^s+o(1)\quad\Longleftrightarrow\quad
 c(z-1)=s+o(1).
\]
\item If $c,d\ge1$, $d/c=t+o(1)$ with $t\in\R_{>0}$, and
$z^c=r+o(1)$ with $r>0$ real, then $z^d=r^t+o(1)$.
\item If $c$ is finite and $\alpha=\st(c)$, then
$z^c=r+o(1)$, $r>0$, implies $z=r^{1/\alpha}+o(1)$.
\end{enumerate}
\end{lemma}

\begin{lemma}\label{sf:lem:atoms}
Every member of $\Sk$ is a finite nonempty sum of additive irreducibles.
Every additive irreducible other than $1$ and $x$ admits one of the forms
\[
 uv\quad(u,v\in\Sk,\ u,v\ge x),\qquad
 u^v\quad(u,v\in\Sk,\ u\ge2,\ v\ge x).
\]
\end{lemma}

\begin{lemma}\label{sf:lem:subsequence}
Every sequence in a well-order has an infinite weakly increasing
subsequence. The same is true simultaneously for finitely many
sequences in well-orders.
\end{lemma}

\begin{theorem}[Berarducci--Mamino discreteness in $F$]\label{sf:thm:discrete}
For every $Q\in\Sk$ and every $c\in F$ with $c\ge1$, the set
\[
 D_c(Q)=\{r\in\R_{>0}:(f/Q)^c=r+o(1)
                         \text{ for some }f\in\Sk\}
\]
is well ordered and has no accumulation point in $\R$.
In particular $D_c(Q)\cap(0,R]$ is finite for every real $R>0$,
and $|D_c(Q)|\le\omega$.
\end{theorem}

\begin{proof}
The map $f\mapsto\st((f/Q)^c)$ on those $f$ with positive finite
standard part is weakly increasing. Its image is therefore well
ordered: for a nonempty subset of the image, choose the least
element of its inverse image in $\Sk$.

A well-ordered subset of $\R_{>0}$ has a positive minimum, and so
cannot accumulate at zero. If it has an accumulation point anywhere
else, it has a strictly increasing bounded sequence. Indeed an
infinite bounded subset has order type at least $\omega$; its first
$\omega$ elements give such a sequence. It suffices to rule these out.

Assume a counterexample exists. Choose the least $Q\in\Sk$ for
which a counterexample exists for \emph{some} $c\in F_{\ge1}$.
Thus, for every $P\in\Sk$ with $P<Q$ and every $d\in F_{\ge1}$,
a bounded weakly increasing sequence of positive elements of $D_d(P)$
is eventually constant. We refer to this as the reference induction.

First we record a consequence which does not use any further
minimality: no counterexample at $Q$ can consist entirely of additive
irreducibles. Suppose otherwise, writing
\[
 (a_i/Q)^c=r_i+o(1),\qquad r_0<r_1<\cdots<R.
\]
Then $a_i$ is strictly increasing and $a_i\arch Q$.
The exceptions $1,x$ can occur at most twice. By
Lemma~\ref{sf:lem:atoms} and a subsequence, all remaining $a_i$ have
one of its two forms, with both coordinate sequences weakly
increasing by Lemma~\ref{sf:lem:subsequence}. Replacing the reference
in the displayed ratios by $a_0$ divides every $r_i$ by $r_0$.

If $a_i=u_iv_i$ with $u_i,v_i\ge x$, then
\[
 (a_i/a_0)^c=(u_i/u_0)^c(v_i/v_0)^c.
\]
Both factors are at least $1$ and have bounded positive real standard
parts. Their standard parts are weakly increasing. As
$u_0,v_0\prec a_0\arch Q$, reference induction makes both
sequences eventually constant, contradicting strict increase of $r_i$.

If $a_i=u_i^{v_i}$ with $u_i\ge2$ and $v_i\ge x$, then
\[
 (a_i/a_0)^c\ge a_i/a_0\ge2^{v_i-v_0}.
\]
The left side has standard part uniformly bounded in $i$; enlarging
that real bound by $1$ gives a uniform bound in $F$ as well. Hence
$v_i-v_0$ is bounded by one integer. Lemma~\ref{sf:lem:integers}
and monotonicity show that $v_i$ is eventually constant. After
discarding an initial segment and renumbering, all $v_i=v_0$ and
\[
 (a_i/a_0)^c=(u_i/u_0)^{v_0c}.
\]
Here $u_0\prec Q$, and $v_0c\in F_{\ge1}$. Reference induction
again contradicts strict increase. This proves the assertion about
additive irreducibles for every exponent and every counterexample at $Q$.

We now choose a counterexample of the form
\begin{equation}
 (h_i/Q)^c=r_i+o(1),\qquad 0<r_0<r_1<\cdots<R,
 \label{sf:eq:witness}
\end{equation}
with an additional minimality condition. Let $A$ be the least Skolem
function Archimedean-equivalent to $Q$. For every counterexample
at $Q$, all its members satisfy $h_i\le NA$ for some common
positive integer $N$. To see this, if $h_i/Q\ge1$, then
$h_i/Q\le(h_i/Q)^c<R+1$; if $h_i/Q<1$ there is already such
a bound. Combine this with $Q=O(A)$. Choose the smallest $N$
among all counterexamples at $Q$, allowing the exponent $c$ to vary.
Fix one such witness. Passing to subsequences preserves its bound.

The element $Q$ is the least Skolem member of its $\arch_c$-class.
Otherwise a smaller reference $Q'\arch_c Q$ would multiply the
real sequence in \eqref{sf:eq:witness} by a fixed positive real constant,
contradicting minimality of $Q$.

Only finitely many $h_i$ can be additive irreducibles, by the assertion
already proved. Decompose every remaining $h_i$ into additive irreducibles,
choose one of maximal value and call it $f_i$, and let $g_i$ be the
sum of the others. Then
\begin{equation}
 h_i=f_i+g_i,\quad f_i,g_i\in\Sk,\quad
 f_i\text{ an additive irreducible},\quad f_i\arch h_i\arch Q.
 \label{sf:eq:split}
\end{equation}
The comparison $f_i\arch h_i$ follows because the decomposition
has a finite number of summands for each $i$; no uniform bound on
that number is used. Pass to a subsequence on which both $f_i$ and
$g_i$ are weakly increasing. In particular $f_i\ge A$.

If $c$ is finite, Lemma~\ref{sf:lem:scaled}(iv) changes the witness
to one with exponent $1$, preserving strict increase, boundedness,
and the integer bound $N$. Minimality of $Q$ now implies $A=Q$.
Let
\[
 a_i=\st(f_i/Q)>0,\qquad b_i=\st(g_i/Q)\ge0.
\]
Both sequences are weakly increasing and bounded, and their sum is
the strictly increasing ratio sequence. If all $b_i=0$, the
sequence $f_i$ contradicts the assertion about additive irreducibles.
Otherwise discard an initial segment so that every $b_i>0$.
Then $f_i,g_i\arch Q$, so both are at least $A=Q$, and
\[
 f_i,g_i\le(N-1)A.
\]
If either $a_i$ or $b_i$ were not eventually constant, a strictly
increasing subsequence would be a counterexample at $Q$ with smaller
integer bound. Therefore both are eventually constant, a contradiction.

It remains to treat infinite $c$. Normalize at $h_0$ and put
\[
 (h_i/h_0)^c=R_i+o(1),\qquad s_i=\log R_i.
\]
Then $R_i=r_i/r_0$, and $s_i$ is a strictly increasing bounded
sequence of nonnegative reals. Lemma~\ref{sf:lem:scaled}(ii), followed
by \eqref{sf:eq:split}, gives
\[
 c(f_i/f_0-1)+(c/f_0)(g_i-g_0)
       =s_i(1+g_0/f_0)+o(1).
\]
Both summands on the left are nonnegative and finite. Define
\begin{equation}
 a_i=\st\bigl(c(f_i/f_0-1)\bigr),\qquad
 b_i=\st\bigl((c/f_0)(g_i-g_0)\bigr),\qquad
 q=\st(g_0/f_0)\ge0.
 \label{sf:eq:ab}
\end{equation}
The two sequences are weakly increasing and bounded, and
\begin{equation}
 a_i+b_i=(1+q)s_i.\label{sf:eq:sum-ab}
\end{equation}

The sequence $a_i$ is eventually constant. Indeed
$(f_i/f_0)^c=e^{a_i}+o(1)$. If $f_0<Q$, use reference induction.
If $Q\le f_0$, then $Q\le f_0\le h_0\arch_c Q$ puts $f_0$
in the same $\arch_c$-class by convexity. Thus a nonconstant
$a_i$ would yield a counterexample at $Q$ consisting entirely of
the additive irreducibles $f_i$, which has been excluded.

Set $P=f_0/c$. Then $P\prec Q$, and the definition of $b_i$ gives
\begin{equation}
 g_i-g_0=b_iP+o(P).\label{sf:eq:remainder}
\end{equation}
We prove that $b_i$ is eventually constant by three cases.

If $g_0\prec P$, then $g_i/P=b_i+o(1)$. If every $b_i=0$
there is nothing to prove. Otherwise choose $k$ with $b_k>0$.
For $i\ge k$ the actual Skolem reference $g_k$ satisfies
\[
 g_k\sim b_kP\prec Q,\qquad
 g_i/g_k=b_i/b_k+o(1).
\]
Reference induction with exponent $1$ proves the assertion.
If $g_0\arch P$, write $g_0/P=t+o(1)$ with $t>0$ real.
Then $g_0\prec Q$ and
$g_i/g_0=1+b_i/t+o(1)$, so the same induction applies to $g_0$.
In neither case is $P$ used as a Skolem reference.

Finally suppose $g_0\succ P$, and set $d=g_0/P=(g_0/f_0)c$.
This is an infinite element of $F$. Dividing \eqref{sf:eq:remainder}
by $g_0$ and using Lemma~\ref{sf:lem:scaled}(ii) yields
\[
 (g_i/g_0)^d=e^{b_i}+o(1).
\]
If $g_0<Q$, reference induction applies. Otherwise
$Q\le g_0\le h_0\arch_c Q$, so $g_0\arch_c Q$.
Moreover
\[
 Q\le g_0\le h_0\arch Q,\qquad f_0\arch h_0\arch Q,
\]
so $g_0$ and $f_0$ are both Archimedean-equivalent to $Q$.
Thus $g_0\arch f_0$, and $c/d=\gamma+o(1)$ for some
$\gamma\in\R_{>0}$. Lemma~\ref{sf:lem:scaled}(iii) gives
\[
 (g_i/Q)^c=R_*e^{\gamma b_i}+o(1)
\]
for a fixed $R_*>0$ real. As $f_i\ge A$ and $h_i\le NA$,
we have $g_i\le(N-1)A$. A nonconstant $b_i$ would therefore
contradict the choice of $N$. This completes all three cases.

Both sequences on the left of \eqref{sf:eq:sum-ab} are eventually
constant, whereas its right side is strictly increasing. The
contradiction proves the theorem. Finiteness in each bounded real
interval follows from compactness, including the exclusion of
accumulation at zero, and implies the order-type assertion.
\end{proof}

\end{document}
